\documentclass[UTF8]{ctexart}
\usepackage{amsmath, amssymb, geometry}
\usepackage{xcolor}
\usepackage{enumitem}
\usepackage{tcolorbox}
\usepackage{cases}
\usepackage{fancyhdr}
\geometry{a4paper, margin=1in}

% 定义红色环境
\newenvironment{redenv}{\color{red}}{}

% 定义详细解析环境
\newenvironment{detailedanalysis}{%
    \color{red}
    \begin{enumerate}[label=\textbf{第\arabic*步：}, leftmargin=*, align=left, itemsep=1.5ex]
}{%
    \end{enumerate}
}

% 定义概念框
\newenvironment{conceptbox}{%
    \begin{tcolorbox}[colback=red!5, colframe=red!50!black, title=概念解析, fonttitle=\bfseries]
}{%
    \end{tcolorbox}
}

% 定义公式推导环境
\newenvironment{formuladerivation}{%
    \begin{enumerate}[label={\textbf{公式\arabic*：}}, leftmargin=*, align=left]
}{%
    \end{enumerate}
}

% 定义题目和解析的分隔线
\newcommand{\problemrule}{\noindent\rule{\textwidth}{1.5pt}\vspace{-0.8em}}
\newcommand{\analysisrule}{\noindent\rule{\textwidth}{0.8pt}\vspace{-0.8em}}

% 宏定义：方便排版选择题选项
\newcommand{\choicesFour}[4]{
    \begin{enumerate}[label=\Alph*., itemsep=0pt, parsep=0pt, topsep=5pt, labelsep=0.5em]
        \begin{minipage}{0.22\linewidth} \item #1 \end{minipage}
        \begin{minipage}{0.22\linewidth} \item #2 \end{minipage}
        \begin{minipage}{0.22\linewidth} \item #3 \end{minipage}
        \begin{minipage}{0.22\linewidth} \item #4 \end{minipage}
    \end{enumerate}
}
\newcommand{\choicesTwo}[4]{
    \begin{enumerate}[label=\Alph*., itemsep=0pt, parsep=0pt, topsep=5pt, labelsep=0.5em]
        \begin{minipage}{0.45\linewidth} \item #1 \end{minipage}
        \begin{minipage}{0.45\linewidth} \item #2 \end{minipage}
        \begin{minipage}{0.45\linewidth} \item #3 \end{minipage}
        \begin{minipage}{0.45\linewidth} \item #4 \end{minipage}
    \end{enumerate}
}
\newcommand{\choices}[4]{
    \begin{enumerate}[label=\Alph*., itemsep=0pt, parsep=0pt, topsep=5pt, labelsep=0.5em]
        \item #1
        \item #2
        \item #3
        \item #4
    \end{enumerate}
}

\newcommand{\blank}[1]{\underline{\hspace{#1}}}

\begin{document}

\section*{第一章 函数（提高）}

% ========== 第1题 ==========
\problemrule
\section*{【题目1】复合函数求值}
设 $ f(x)=\sin x $ , $ g(x)=\begin{cases}x-\pi, & x\leq0,\\x+\pi, & x>0,\end{cases} $ 则 $ f[g(x)]= $ \blank{1cm}。
\choicesFour{$ \sin x $}{$ \cos x $}{$ - \sin x $}{$ - \cos x $}

\begin{redenv}
\analysisrule
\subsection*{逐步解析}

\begin{detailedanalysis}
\item \textbf{问题本质分析}
本题考查分段函数与三角函数的复合运算。核心在于根据外层函数 $f(x)$ 的定义，处理内层函数 $g(x)$ 的不同分段情况。

\item \textbf{解题策略构建}
\begin{enumerate}[label=(\alph*)]
\item 分析内层函数 $g(x)$ 的分段条件。
\item 分别讨论 $x \le 0$ 和 $x > 0$ 时，$g(x)$ 的表达式。
\item 代入外层函数 $f(u) = \sin u$，利用诱导公式化简。
\item 观察两段结果是否一致，从而合并。
\end{enumerate}

\item \textbf{详细求解过程}
\textbf{情形1：当 $x \le 0$ 时}
$$ g(x) = x - \pi $$
$$ f[g(x)] = f(x - \pi) = \sin(x - \pi) $$
利用诱导公式 $\sin(\alpha - \pi) = -\sin(\pi - \alpha) = -\sin \alpha$（或直接 $\sin(x-\pi) = -\sin x$）。
$$ f[g(x)] = -\sin x $$

\textbf{情形2：当 $x > 0$ 时}
$$ g(x) = x + \pi $$
$$ f[g(x)] = f(x + \pi) = \sin(x + \pi) $$
利用诱导公式 $\sin(x + \pi) = -\sin x$。
$$ f[g(x)] = -\sin x $$

\textbf{综合}
无论 $x$ 取何值，都有 $f[g(x)] = -\sin x$。

\item \textbf{最终答案}
\textbf{答案}：C
\end{detailedanalysis}

\subsection*{核心公式与推导}
\begin{formuladerivation}
\item \textbf{三角函数诱导公式}
$$ \sin(\pi + \alpha) = -\sin \alpha $$
$$ \sin(\alpha - \pi) = -\sin(\pi - \alpha) = -\sin \alpha $$
推导：在单位圆中，角 $\alpha$ 与 $\alpha + \pi$ 的终边关于原点对称，正弦值互为相反数。
\end{formuladerivation}

\subsection*{考点深度分析}
\begin{conceptbox}
\textbf{知识点归类}：复合函数、分段函数、三角函数性质。
\end{conceptbox}
\noindent\textbf{易错点}：诱导公式符号记忆错误，或者分别求出分段结果后忘记观察是否可以合并。
\end{redenv}

% ========== 第2题 ==========
\problemrule
\section*{【题目2】具体函数定义域}
函数 $ y = \ln | \sin x | $ 的定义域是 \blank{1cm}，其中 $k$ 为整数。
\choicesTwo
{$ x \neq \frac{k\pi}{2} $}
{$ x \in (-\infty, \infty), x \neq k\pi $}
{$ x = k\pi $}
{$ x \in (-\infty, \infty) $}

\begin{redenv}
\analysisrule
\subsection*{逐步解析}

\begin{detailedanalysis}
\item \textbf{问题本质分析}
求复合函数的定义域，主要考虑对数函数和三角函数的限制条件。

\item \textbf{详细求解过程}
\begin{enumerate}[label=(\arabic*)]
\item \textbf{对数限制}：$\ln u$ 要求真数 $u > 0$。即 $|\sin x| > 0$。
\item \textbf{绝对值性质}：$|a| \ge 0$，且 $|a|=0$ 当且仅当 $a=0$。
\item \textbf{求解不等式}：要使 $|\sin x| > 0$，只需 $\sin x \neq 0$。
\item \textbf{三角方程求解}：$\sin x = 0$ 的解为 $x = k\pi$ ($k \in \mathbb{Z}$)。
\item \textbf{结论}：定义域为 $x \neq k\pi$。
\end{enumerate}

\item \textbf{最终答案}
\textbf{答案}：B
\end{detailedanalysis}

\subsection*{考点深度分析}
\begin{conceptbox}
\textbf{知识点归类}：函数定义域。
\end{conceptbox}
\noindent\textbf{易错点}：忽略绝对值的作用，误以为 $\sin x > 0$。
\end{redenv}

% ========== 第3题 ==========
\problemrule
\section*{【题目3】复合函数定义域}
函数 $ f(x)=\frac{1-x^{2}}{\sqrt{|x|-1}}+\arcsin(x-1) $ 的定义域是 \blank{1cm}。
\choicesFour{$ (0,2] $}{$ [0,2] $}{$ (1,2] $}{$ [1,2] $}

\begin{redenv}
\analysisrule
\subsection*{逐步解析}

\begin{detailedanalysis}
\item \textbf{问题本质分析}
本题涉及分式、根式、反三角函数三部分的定义域交集。

\item \textbf{详细求解过程}
需要满足以下三个条件：
\begin{enumerate}[label=(\arabic*)]
\item \textbf{分母不为零且根号内大于0（针对分母位置的偶次根式）}：
$$ |x| - 1 > 0 \implies |x| > 1 \implies x > 1 \text{ 或 } x < -1 $$
\item \textbf{反正弦函数定义域}：
$$ -1 \le x-1 \le 1 $$
$$ 0 \le x \le 2 $$
\item \textbf{求交集}：
$$ \{ x | x > 1 \text{ 或 } x < -1 \} \cap \{ x | 0 \le x \le 2 \} $$
在数轴上取交集：
$$ (1, 2] $$
\end{enumerate}

\item \textbf{最终答案}
\textbf{答案}：C
\end{detailedanalysis}

\subsection*{核心公式与推导}
\begin{formuladerivation}
\item \textbf{基本初等函数定义域}
$$ y = \sqrt{x} \implies x \ge 0 $$
$$ y = \frac{1}{x} \implies x \neq 0 $$
$$ y = \arcsin x \implies -1 \le x \le 1 $$
\end{formuladerivation}
\end{redenv}

% ========== 第4题 ==========
\problemrule
\section*{【题目4】复杂函数定义域}
函数 $ y=\sqrt{2x-x^{2}}-\arcsin\frac{2x-1}{7} $ 的定义域为 \blank{1cm}。
\choicesFour{$ \left[-3,4\right] $}{$ (-3,4) $}{$ [0,2] $}{$ (0,2) $}

\begin{redenv}
\analysisrule
\subsection*{逐步解析}

\begin{detailedanalysis}
\item \textbf{详细求解过程}
需满足：
\begin{enumerate}[label=(\arabic*)]
\item \textbf{偶次根式}：
$$ 2x - x^2 \ge 0 \implies x(2-x) \ge 0 \implies 0 \le x \le 2 $$
\item \textbf{反三角函数}：
$$ -1 \le \frac{2x-1}{7} \le 1 $$
$$ -7 \le 2x - 1 \le 7 $$
$$ -6 \le 2x \le 8 $$
$$ -3 \le x \le 4 $$
\item \textbf{求交集}：
$$ [0, 2] \cap [-3, 4] = [0, 2] $$
\end{enumerate}

\item \textbf{最终答案}
\textbf{答案}：C
\end{detailedanalysis}
\end{redenv}

% ========== 第5题 ==========
\problemrule
\section*{【题目5】对数与反三角定义域}
函数 $ y = \arcsin(1-x) + \frac{1}{2}\lg\frac{1+x}{1-x} $ 的定义域是 \blank{1cm}。
\choicesFour{$ (0,1) $}{$ [0,1) $}{$ [0,1] $}{$ [0,1] $}

\begin{redenv}
\analysisrule
\subsection*{逐步解析}

\begin{detailedanalysis}
\item \textbf{详细求解过程}
\begin{enumerate}[label=(\arabic*)]
\item \textbf{反三角函数}：
$$ -1 \le 1-x \le 1 $$
$$ -2 \le -x \le 0 \implies 0 \le x \le 2 $$
\item \textbf{对数真数大于0}：
$$ \frac{1+x}{1-x} > 0 $$
等价于 $(1+x)(1-x) > 0$，即 $(x+1)(x-1) < 0$。
解得：$-1 < x < 1$。
\item \textbf{求交集}：
$$ [0, 2] \cap (-1, 1) = [0, 1) $$
\end{enumerate}

\item \textbf{最终答案}
\textbf{答案}：B
\end{detailedanalysis}
\end{redenv}

% ========== 第6题 ==========
\problemrule
\section*{【题目6】复合函数求值}
已知函数 $ f(x)=x^{3}+3x-2 $ , $ g(x)=\tan x $ , 则 $ f[g(\frac{\pi}{4})]= $ \blank{2cm}。

\begin{redenv}
\analysisrule
\subsection*{逐步解析}

\begin{detailedanalysis}
\item \textbf{详细求解过程}
\begin{enumerate}[label=(\arabic*)]
\item 先求内层函数值：
$$ g(\frac{\pi}{4}) = \tan \frac{\pi}{4} = 1 $$
\item 再求外层函数值：
$$ f[g(\frac{\pi}{4})] = f(1) $$
代入 $f(x)$ 表达式：
$$ f(1) = 1^3 + 3 \times 1 - 2 = 1 + 3 - 2 = 2 $$
\end{enumerate}

\item \textbf{最终答案}
\textbf{答案}：2
\end{detailedanalysis}
\end{redenv}

% ========== 第7题 ==========
\problemrule
\section*{【题目7】分段复合求值}
已知 $ f(x)=\begin{cases}\frac{1}{x}, & |x|>1\\0, & |x|\leq1\end{cases} $ 则 $ f[f(2021)]= $ \blank{2cm}。

\begin{redenv}
\analysisrule
\subsection*{逐步解析}

\begin{detailedanalysis}
\item \textbf{详细求解过程}
\begin{enumerate}[label=(\arabic*)]
\item \textbf{第一层}：求 $f(2021)$。
因为 $2021 > 1$，满足 $|x|>1$，故
$$ f(2021) = \frac{1}{2021} $$
\item \textbf{第二层}：求 $ f[f(2021)] = f(\frac{1}{2021}) $。
设 $x_2 = \frac{1}{2021}$。
显然 $0 < x_2 < 1$，即 $|x_2| < 1$。
根据分段函数定义，当 $|x| \le 1$ 时，$f(x) = 0$。
所以 $ f(\frac{1}{2021}) = 0 $。
\end{enumerate}

\item \textbf{最终答案}
\textbf{答案}：0
\end{detailedanalysis}
\end{redenv}

% ========== 第8题 ==========
\problemrule
\section*{【题目8】多层对数定义域}
函数 $ y = \ln[\ln(\ln x)] $ 的定义域是 \blank{2cm}。

\begin{redenv}
\analysisrule
\subsection*{逐步解析}

\begin{detailedanalysis}
\item \textbf{详细求解过程}
要使函数有意义，需满足由内而外的限制条件：
\begin{enumerate}[label=(\arabic*)]
\item 最内层对数 $\ln x$：$x > 0$。
\item 中间层对数 $\ln(\ln x)$：要求真数 $\ln x > 0 \implies x > 1$。
\item 最外层对数 $\ln[\ln(\ln x)]$：要求真数 $\ln(\ln x) > 0$。
$$ \ln(\ln x) > 0 \implies \ln x > 1 \implies x > e $$
\item \textbf{求交集}：
同时满足 $x>0, x>1, x>e$，取最强条件。
$$ \therefore x > e $$
\end{enumerate}

\item \textbf{最终答案}
\textbf{答案}：$ (e, +\infty) $
\end{detailedanalysis}
\end{redenv}

% ========== 第9题 ==========
\problemrule
\section*{【题目9】分式根式定义域}
函数 $ f(x)=\frac{1}{\sqrt{x-3}} $ 的定义域是 \blank{2cm}。

\begin{redenv}
\analysisrule
\subsection*{逐步解析}

\begin{detailedanalysis}
\item \textbf{详细求解过程}
分母中的偶次根式需满足：被开方数大于0（因为分母不能为0）。
$$ x - 3 > 0 \implies x > 3 $$

\item \textbf{最终答案}
\textbf{答案}：$ (3, +\infty) $
\end{detailedanalysis}
\end{redenv}

% ========== 第10题 ==========
\problemrule
\section*{【题目10】综合定义域}
函数 $ f(x) = \sqrt{4 - x^{2}} + \frac{1}{\ln \cos x} $ 的定义域为 \blank{2cm}。

\begin{redenv}
\analysisrule
\subsection*{逐步解析}

\begin{detailedanalysis}
\item \textbf{详细求解过程}
\begin{enumerate}[label=(\arabic*)]
\item \textbf{根式}：$4 - x^2 \ge 0 \implies x^2 \le 4 \implies -2 \le x \le 2$。
\item \textbf{对数分母}：
\begin{itemize}
\item 对数真数大于0：$\cos x > 0$。在 $[-2, 2]$ 范围内（注意 $2 \approx 0.63\pi$），$\cos x > 0 \implies x \in (-\frac{\pi}{2}, \frac{\pi}{2})$。
\item 分母不为零：$\ln \cos x \neq 0 \implies \cos x \neq 1 \implies x \neq 0$。
\end{itemize}
\item \textbf{求交集}：
区间 $[-2, 2]$ 与 $(-\frac{\pi}{2}, \frac{\pi}{2})$。因为 $\frac{\pi}{2} \approx 1.57 < 2$，所以交集受限于 $\frac{\pi}{2}$。
但是题目限制了 $x \in [-2, 2]$。不，$\cos x > 0$ 是必要条件。
$x \in [-2, 2]$ 且 $-\frac{\pi}{2} < x < \frac{\pi}{2}$ 且 $x \neq 0$。
由于 $[-2, 2]$ 覆盖了 $(-\frac{\pi}{2}, \frac{\pi}{2})$（不对，$2 > 1.57$），所以限制实际上是 $(-\frac{\pi}{2}, \frac{\pi}{2})$ 去掉 $0$。
\end{enumerate}

\item \textbf{最终答案}
\textbf{答案}：$ (-\frac{\pi}{2}, 0) \cup (0, \frac{\pi}{2}) $
\end{detailedanalysis}
\end{redenv}

% ========== 第11题 ==========
\problemrule
\section*{【题目11】函数解析式}
设 $ f(x-1)=x(x-1) $ ，则 $ f(x)= $ \blank{2cm}。

\begin{redenv}
\analysisrule
\subsection*{逐步解析}

\begin{detailedanalysis}
\item \textbf{解题策略}
换元法或配凑法。

\item \textbf{详细求解过程}
\textbf{方法一：配凑法}
观察右边 $x(x-1)$，试图配出 $(x-1)$。
$$ f(x-1) = x(x-1) = (x-1+1)(x-1) = (x-1)^2 + (x-1) $$
令 $t = x-1$，则 $f(t) = t^2 + t$。
所以 $f(x) = x^2 + x$。

\textbf{方法二：换元法}
令 $t = x-1 \implies x = t+1$。
代入原式：
$$ f(t) = (t+1)t = t^2 + t $$
替换变量名为 $x$：
$$ f(x) = x^2 + x $$

\item \textbf{最终答案}
\textbf{答案}：$ x^2 + x $
\end{detailedanalysis}
\end{redenv}

% ========== 第12题 ==========
\problemrule
\section*{【题目12】偶函数延拓}
设 $ f(x) $ 是 $ (-\infty, +\infty) $ 内的偶函数，并且当 $ x \in (-\infty, 0) $ 时，有 $ f(x) = x + 2 $ ，则当 $ x \in (0, +\infty) $ 时， $ f(x) $ 的表达式是 \blank{2cm}。

\begin{redenv}
\analysisrule
\subsection*{逐步解析}

\begin{detailedanalysis}
\item \textbf{解题策略}
利用偶函数定义 $f(x) = f(-x)$。

\item \textbf{详细求解过程}
设 $x > 0$，则 $-x < 0$。
因为当 $t < 0$ 时，$f(t) = t + 2$。
所以 $f(-x) = (-x) + 2 = -x + 2$。
又因为 $f(x)$ 是偶函数，所以 $f(x) = f(-x)$。
$$ \therefore f(x) = -x + 2 \quad (x>0) $$

\item \textbf{最终答案}
\textbf{答案}：$ -x+2 $
\end{detailedanalysis}
\end{redenv}

% ========== 第13题 ==========
\problemrule
\section*{【题目13】函数值域}
函数 $ y = \sin x - \sin |x| $ 的值域是 \blank{2cm}。

\begin{redenv}
\analysisrule
\subsection*{逐步解析}

\begin{detailedanalysis}
\item \textbf{解题策略}
去绝对值，分段讨论。

\item \textbf{详细求解过程}
\begin{enumerate}[label=(\arabic*)]
\item 当 $x \ge 0$ 时：$|x| = x$。
$$ y = \sin x - \sin x = 0 $$
\item 当 $x < 0$ 时：$|x| = -x$。
$$ y = \sin x - \sin(-x) = \sin x - (-\sin x) = 2\sin x $$
因为 $x < 0$，$\sin x$ 取值范围为 $[-1, 1]$（注意 $x$ 取负无穷到0）。
更精确地说，$\sin x$ 在 $(-\infty, 0)$ 上可以取到 $[-1, 1]$ 之间的所有值。
所以 $2\sin x$ 的范围是 $[-2, 2]$。
\item \textbf{综合}：
值域是 $\{0\} \cup [-2, 2] = [-2, 2]$。
\end{enumerate}

\item \textbf{最终答案}
\textbf{答案}：$ [-2, 2] $
\end{detailedanalysis}
\end{redenv}

% ========== 第14题 ==========
\problemrule
\section*{【题目14】复合函数求值}
设函数 $ f(x)=x^{3}-x^{2}-1 $ , 则 $ f[f(1)]= $ \blank{2cm}。

\begin{redenv}
\analysisrule
\subsection*{逐步解析}

\begin{detailedanalysis}
\item \textbf{详细求解过程}
\begin{enumerate}[label=(\arabic*)]
\item \textbf{求内层 $f(1)$}：
$$ f(1) = 1^3 - 1^2 - 1 = 1 - 1 - 1 = -1 $$
\item \textbf{求外层 $f[f(1)] = f(-1)$}：
$$ f(-1) = (-1)^3 - (-1)^2 - 1 = -1 - 1 - 1 = -3 $$
\end{enumerate}

\item \textbf{最终答案}
\textbf{答案}：-3
\end{detailedanalysis}
\end{redenv}

% ========== 第15题 ==========
\problemrule
\section*{【题目15】复合函数表达式}
已知函数 $ f(x)=\frac{x+1}{x-1},x\in(1,+\infty) $ ，求复合函数 $ f[f(x)] $ 。

\begin{redenv}
\analysisrule
\subsection*{逐步解析}

\begin{detailedanalysis}
\item \textbf{详细求解过程}
直接代入表达式计算：
$$ f[f(x)] = \frac{f(x)+1}{f(x)-1} = \frac{\frac{x+1}{x-1}+1}{\frac{x+1}{x-1}-1} $$
分子分母同时乘以 $(x-1)$（注意 $x>1$，所以 $x-1 \neq 0$）：
$$ = \frac{(x+1) + (x-1)}{(x+1) - (x-1)} = \frac{2x}{2} = x $$
注意定义域：内层 $x>1$；外层要求 $f(x)>1$。
$f(x) = \frac{x+1}{x-1} = 1 + \frac{2}{x-1}$。当 $x>1$ 时，$\frac{2}{x-1}>0$，所以 $f(x)>1$ 恒成立。
所以 $f[f(x)] = x$。

\item \textbf{最终答案}
\textbf{答案}：$ x $
\end{detailedanalysis}
\end{redenv}

\end{document}
